% ============================================================================
%  atlas-preamble.tex — standalone AI Chip Fab Atlas publication style
%  Engine: pdfLaTeX. Compile from report/ with: latexmk -pdf atlas.tex
% ============================================================================

% ---------------------------------------------------------------- encoding --
\usepackage[T1]{fontenc}
\usepackage[utf8]{inputenc}
\usepackage[english]{babel}
\usepackage{csquotes}

% ------------------------------------------------------------------- fonts --
% XCharter (Bitstream Charter, extended): sturdy, economical, editorial.
% `sups' gives true superior figures for footnote markers.
\usepackage[scaled=1.04,sups]{XCharter}
% Matching math. Drop `vvarbb' if you prefer the standard blackboard bold.
\usepackage[charter,vvarbb,scaled=1.05]{newtxmath}
\usepackage[semibold]{sourcesanspro}          % sans: headings, figures, tables
\usepackage[scaled=0.94,varqu,varl]{inconsolata} % mono
% Text-layer mappings for ligature glyphs. pdfTeX's standard table lacks the
% underscore names Source Sans Pro uses (f_f, f_t, f_f_t), so without these
% lines words like "Office" and "after" copy and search as "Oice" and "aer".
\input{glyphtounicode}
\pdfglyphtounicode{f_f}{0066 0066}
\pdfglyphtounicode{f_i}{0066 0069}
\pdfglyphtounicode{f_l}{0066 006C}
\pdfglyphtounicode{f_f_i}{0066 0066 0069}
\pdfglyphtounicode{f_f_l}{0066 0066 006C}
\pdfglyphtounicode{f_t}{0066 0074}
\pdfglyphtounicode{f_f_t}{0066 0066 0074}
\pdfglyphtounicode{f_j}{0066 006A}
\pdfglyphtounicode{f_b}{0066 0062}
\pdfglyphtounicode{f_h}{0066 0068}
\pdfglyphtounicode{f_k}{0066 006B}
\pdfgentounicode=1

\linespread{1.06}   % XCharter has a large x-height; it wants a little more air

% ------------------------------------------------------------------ colours --
\usepackage[dvipsnames,table]{xcolor}
\definecolor{accent}{HTML}{33393F}   % graphite — headings, rules, links
\definecolor{accentlt}{HTML}{5A6470}
\definecolor{hairline}{HTML}{C9CBCE}
\definecolor{softbg}{HTML}{F3F4F5}
\definecolor{inktext}{HTML}{1A1A1A}  % not pure black; easier on the eye
\color{inktext}

% ---------------------------------------------------------------- geometry --
% Symmetric text block. 5.90in at 12pt gives ~79 characters per line.
% Widen further only if you also raise the point size, or the measure gets long.
\KOMAoptions{fontsize=12pt}   % set here so the style is self-contained
\usepackage[letterpaper]{geometry}
\geometry{
  textwidth  = 5.90in,
  textheight = 8.60in,
  centering,
  headsep    = 0.26in,
  footskip   = 0.55in,
}

% ------------------------------------------------------- microtypography ----
\usepackage{microtype}
\microtypesetup{protrusion=true,expansion=true,tracking=true,factor=1100}
\SetTracking{encoding={*},shape=sc}{45}   % letterspace small caps

% ------------------------------------------------------ page discipline -----
\widowpenalty=10000
\clubpenalty=10000
\brokenpenalty=10000
\raggedbottom                 % prefer short pages over stretched leading
\setlength{\emergencystretch}{2em}
\setlength{\parindent}{1.2em}

% ------------------------------------------------------------- headings -----
\setkomafont{disposition}{\normalcolor\sffamily}
\setkomafont{section}{\sffamily\bfseries\large\color{accent}}
\setkomafont{subsection}{\sffamily\bfseries\normalsize\color{accent}}
\setkomafont{subsubsection}{\sffamily\itshape\normalsize\color{accent}}
\setkomafont{paragraph}{\sffamily\bfseries\small}
\setkomafont{descriptionlabel}{\sffamily\bfseries}

\RedeclareSectionCommand[beforeskip=-3.2ex plus -1ex minus -.2ex,
                         afterskip=1.4ex plus .2ex]{section}
\RedeclareSectionCommand[beforeskip=-2.6ex plus -1ex minus -.2ex,
                         afterskip=1.0ex plus .2ex]{subsection}

% ------------------------------------------------------ heads and feet ------
\usepackage{scrlayer-scrpage}
\pagestyle{scrheadings}
\clearpairofpagestyles
\setkomafont{pagehead}{\sffamily\footnotesize\color{accent!75!black}}
\setkomafont{pagefoot}{\sffamily\footnotesize\color{black!55}}
\automark[section]{section}
\ihead{\runningtitle}
\ohead{\headmark}
\ofoot*{\thepage}
\setlength{\headheight}{14pt}
\KOMAoption{headsepline}{0.4pt}
\addtokomafont{headsepline}{\color{hairline}}

% ------------------------------------------------------------- captions -----
\usepackage[font=small,
            labelfont={sf,bf,color=accent},
            textfont={sf},
            labelsep=period,
            skip=7pt,
            justification=raggedright,
            singlelinecheck=false]{caption}

% ------------------------------------------------------- lists and tables ---
\usepackage{enumitem}
\setlist{itemsep=0.25em, parsep=0pt, topsep=0.5em}
\setlist[itemize,1]{label=\textcolor{accent}{\textbullet}}
\setlist[enumerate,1]{label=\textcolor{accent}{\arabic*.}}

\usepackage{booktabs, tabularx, multirow, siunitx}
\renewcommand{\arraystretch}{1.18}
\newcommand{\tablefont}{\sffamily\small}

% -------------------------------------------------------------- figures -----
\usepackage{graphicx}
\usepackage{float}
\usepackage{pgfplots}
\pgfplotsset{compat=1.17}   % raise once everyone is on TeX Live 2023+
\pgfplotsset{
  policyplot/.style={
    width=\linewidth, height=0.58\linewidth,
    tick label style={font=\sffamily\footnotesize, color=black!60},
    label style={font=\sffamily\small, color=black!75},
    legend style={font=\sffamily\footnotesize, draw=none, fill=none,
                  at={(0.02,0.98)}, anchor=north west},
    title style={font=\sffamily\bfseries\small, color=accent},
    axis line style={color=black!35},
    grid=major, grid style={hairline, line width=0.3pt},
    axis x line*=bottom, axis y line*=left,
    tick align=outside, tick style={color=black!35},
    every axis plot/.append style={line width=1.1pt},
    cycle list={{accent},{Sepia},{accentlt},{OliveGreen},{black!55}},
  }
}

% ------------------------------------------------------------ side notes ----
% No wide margin any more, so \sidenote degrades to a footnote. This keeps any
% existing \sidenote markup working. Delete if you'd rather it error loudly.
\newcommand{\sidenote}[1]{\footnote{#1}}

% --------------------------------------------------------- callout boxes ----
\usepackage[most]{tcolorbox}

% tcolorbox parks \footnote inside the box. Capture them instead and replay
% them once the box has ended, so they land at the foot of the page.
\makeatletter
\newcommand{\tcb@fnstore}{}
\AtBeginDocument{\global\let\tcb@stdfnmark\@makefnmark
  \global\let\tcb@stdmpfn\@mpfn \global\let\tcb@stdthefn\thefootnote}
\newcommand{\tcbfncollect}{%
  % Reset globally: the store is filled via \g@addto@macro, and a local reset
  % made here would be undone when the box's group ends, before the flush.
  \gdef\tcb@fnstore{}%
  % tcolorbox sets up a minipage-like footnote context, which suppresses
  % XCharter's superior figures. Restore the outer context for the marker.
  \let\@makefnmark\tcb@stdfnmark
  \let\@mpfn\tcb@stdmpfn \let\thefootnote\tcb@stdthefn
  \renewcommand{\footnote}[1]{%
    \footnotemark
    \begingroup
      \edef\tcb@fnx{\noexpand\g@addto@macro\noexpand\tcb@fnstore{%
        \noexpand\footnotetext[\the\value{footnote}]{\unexpanded{##1}}}}%
    \expandafter\endgroup\tcb@fnx
  }%
}
% Replay, then clear globally so a later box cannot re-emit these footnotes.
% (\footnote needs no restore: the box-local redefinition dies with the box.)
\newcommand{\tcbfnflush}{\tcb@fnstore\gdef\tcb@fnstore{}}
\makeatother
\newtcolorbox{keybox}[1][Key findings]{
  enhanced, breakable,
  colback=softbg, colframe=accent,
  boxrule=0pt, leftrule=2.5pt,
  arc=0pt, outer arc=0pt,
  left=12pt, right=12pt, top=10pt, bottom=10pt,
  fonttitle=\sffamily\bfseries\small\color{accent},
  title={#1}, coltitle=accent,
  attach boxed title to top left={yshift=-2pt},
  boxed title style={empty},
  before upper={\tcbfncollect}, after={\tcbfnflush},
}

% Untitled, same visual language. For a passage set apart without a label.
\newtcolorbox{plainbox}{
  enhanced, breakable,
  colback=softbg, colframe=accent,
  boxrule=0pt, leftrule=2.5pt,
  arc=0pt, outer arc=0pt,
  left=12pt, right=12pt, top=10pt, bottom=10pt,
  before upper={\tcbfncollect}, after={\tcbfnflush},
}

% Quieter still: hairline rules above and below, no fill, no left bar.
\newtcolorbox{quietbox}{
  enhanced, breakable,
  colback=white, boxrule=0pt, frame hidden,
  arc=0pt, outer arc=0pt,
  borderline north={0.6pt}{0pt}{hairline},
  borderline south={0.6pt}{0pt}{hairline},
  left=0pt, right=0pt, top=11pt, bottom=11pt,
  before upper={\tcbfncollect}, after={\tcbfnflush},
}

% ------------------------------------------------- abstract / summary -------
\newenvironment{execsummary}
  {\begin{center}\begin{minipage}{0.92\linewidth}
   \setlength{\parindent}{0pt}\small
   {\sffamily\bfseries\color{accent}Executive summary}\par\medskip
   \color{black!85}}
  {\end{minipage}\end{center}\bigskip}

% ------------------------------------------------------------ epigraphs -----
\newcommand{\pullquote}[2]{%
  \begin{center}\begin{minipage}{0.85\linewidth}
  \color{accent}\itshape\large #1\par\medskip
  \normalfont\sffamily\footnotesize\color{black!55}\raggedleft #2
  \end{minipage}\end{center}\medskip}

% ---------------------------------------------------------- footnote look ---
\deffootnote{1.2em}{1.2em}{\textsuperscript{\thefootnotemark}\hspace{0.35em}}
\setkomafont{footnote}{\small}

% ------------------------------------------------------------- hyperref -----
\usepackage[unicode]{hyperref}
\hypersetup{
  colorlinks = true,
  linkcolor  = accent,
  citecolor  = accent,
  urlcolor   = accentlt,
  filecolor  = accent,
  breaklinks = true,
  pdfborder  = {0 0 0},
}
\usepackage[nameinlink]{cleveref}

% ---------------------------------------------------------- title matter ----
\newcommand{\runningtitle}{}
\newcommand{\papertitle}{}
\newcommand{\papersubtitle}{}
\newcommand{\paperdate}{\today}
\newcommand{\paperorg}{}

% Author block: each author carries their own affiliation, laid out as a grid
% that wraps automatically. Use \addauthor{Name}{Affiliation} once per author,
% in the preamble of main.tex. For a shared affiliation, just repeat the text —
% or see the grouped alternative in the comments below.
\newcommand{\authorcolwidth}{0.30\linewidth}
\makeatletter
\newcommand{\@authorlist}{}
\newcommand{\addauthor}[2]{%
  \g@addto@macro\@authorlist{%
    \parbox[t]{\authorcolwidth}{\raggedright
      \sffamily\small\bfseries\color{black!85}#1\par
      \vspace{3pt}%
      \normalfont\sffamily\footnotesize\color{black!52}#2\par}%
    \hspace{0.035\linewidth}\allowbreak
  }%
}
\newcommand{\printauthors}{%
  \begingroup
    \setlength{\lineskip}{16pt}%
    \setlength{\parindent}{0pt}%
    \raggedright\@authorlist\par
  \endgroup}
\makeatother

\newcommand{\makecoverpage}{%
  \thispagestyle{empty}
  \begin{flushleft}
    \sffamily
    {\footnotesize\color{black!55}\lsstyle\MakeUppercase{\paperorg}\par}
    \vspace{0.55in}
    {\color{hairline}\rule{\linewidth}{2.5pt}\par}
    \vspace{0.30in}
    % For a serif display title instead, swap \sffamily -> \rmfamily here.
    {\sffamily\bfseries\fontsize{26}{30}\selectfont\color{accent}\papertitle\par}
    \ifx\papersubtitle\empty\else
      \vspace{0.14in}
      {\sffamily\fontsize{15}{19}\selectfont\color{black!62}\papersubtitle\par}
    \fi
    \vspace{0.30in}
    {\color{hairline}\rule{\linewidth}{0.6pt}\par}
    \vspace{0.24in}
    \printauthors
    \vspace{0.22in}
    {\sffamily\footnotesize\color{black!55}\paperdate\par}
  \end{flushleft}
  \vfill
  \newpage
}
